\section{Introduction}
\label{sec:intro}

LLM agents read email, calendars, web pages, and transaction histories on a user's behalf, and every one of those sources can carry text written by someone else. Indirect prompt injection turns that text into instructions~\cite{greshake2023}. The cheapest defense to deploy is a detector: a small classifier that screens every tool output before the agent sees it and drops what looks like an injection~\cite{protectai,piguard,llamafirewall}. It needs no change to the agent or the model, which is why guardrail toolkits ship detectors and practitioners compare them.

The threat is simple to state. An attacker who can write a calendar event description, a product review, or a table cell controls part of what a tool returns; if the detector misses that text, the agent may act on it. If the detector is too eager instead, it drops ordinary tool outputs, and the agent cannot finish the user's task.

How would a team decide which detector to deploy? Detector papers and third-party comparisons report accuracy on prompt collections and injection benchmarks; a 2025 evaluation of nine open guardrails, for example, concluded from BIPIA's email and table tasks that PIGuard detects indirect injection effectively~\cite{mozillaguardrails}. Two things in this practice are unchecked. First, benchmark text is not what an agent's detector screens. Agents see serialized tool outputs (records, lists, confirmations), and false alarms on those translate into failed tasks, not into a percentage on a balanced test set. Second, nobody checks whether the benchmark overlaps with the detector's training data. Prior work showed that same-source splits inflate scores for classifiers the authors train themselves~\cite{benchmarkslie} and argued for low-FPR evaluation~\cite{promptshield}; neither audits released detectors or evaluates them where an agent would use them.

We evaluate fifteen detectors, from 2023 DeBERTa classifiers through Meta's Prompt Guard~2 to 2026 models built for agent inputs, together with two task-aware LLM judges, on the tool outputs of two agent benchmarks, AgentDojo~\cite{agentdojo} and $\tau$-bench~\cite{taubench}, and on BIPIA~\cite{bipia}. Our central observation is that a detector works on inputs whose form resembles what it was trained on, and that benchmark scores mostly measure this resemblance rather than a general ability to detect injections.

Making this measurement requires three things that existing benchmarks do not provide. First, benign agent tool outputs with certain ground truth: we replay the benchmarks' ground-truth tool calls without any LLM, which yields tool outputs that are benign by construction. Second, reliable labels for injected outputs: substring matching fails because AgentDojo re-wraps injected text when it serializes outputs, so we label by differential replay against the clean environment. Third, a way to connect performance to training data: we audit the published training data of two detectors against all three benchmarks and against InjecAgent's attacks~\cite{injecagent}, and we include a judge whose training data predates AgentDojo.

The results reverse the conclusions a benchmark-based comparison would draw. Detection rankings transfer poorly: between BIPIA and AgentDojo Kendall's $\tau$ is \RankTau, and even between the two agent benchmarks it is \TauADTB. PIGuard catches \BIPGTpr of BIPIA injections at 1\% FPR and \ADPGTpr of AgentDojo injections; Prismor catches \MPRAD of AgentDojo injections and \TBTprPR of $\tau$-bench's. False-positive rates on tool outputs, by contrast, range from none to over 90\% and are consistent across the two agent benchmarks ($\tau=\TauFPRADTB$). The training data explains the pattern. PIGuard was trained on \PGBipiaRows full BIPIA inputs; it also saw \NIASeen InjecAgent attacks, but only as short prompts, and inside $\tau$-bench tool outputs it detects those attacks no better than ones it never saw (\TBSeenPG versus \TBUnseenPG). Horizon-Labs shares no data with any benchmark, was trained on agent-style inputs, and leads both agent benchmarks (\MHZAD and \TBTprHZ) with almost no false alarms. Task-aware judges reach 54--78\% on the agent benchmarks, below the best dedicated detectors, and every approach misses whole classes of injections outside its training distribution.

We make the following contributions:
\begin{itemize}
\item A methodology for evaluating injection detectors where agents use them: clean replay of ground-truth tool calls for benign outputs, differential replay for injected ones, and task-level block rates (\S\ref{sec:method}).
\item A measurement of fifteen detectors and two judges on two agent benchmarks and BIPIA, showing that false-positive behavior on tool outputs is stable across agent benchmarks while detection rankings transfer poorly (\S\ref{sec:rq1}, \S\ref{sec:rq2}).
\item Training-data audits and a controlled seen/unseen comparison showing that the form of the training inputs, not overlap in attack strings, decides where a detector works (\S\ref{sec:rq3}).
\item An analysis of what task awareness, attack goals, and input format change, with the finding that task-irrelevant injections evade every approach outside its training distribution (\S\ref{sec:rq4}, \S\ref{sec:rq5}).
\end{itemize}
We will release the code that rebuilds every evaluation set from the public benchmarks and reproduces every number in this paper.
